\documentclass[12pt]{article}
\usepackage[utf8]{inputenc}
\usepackage[T1]{fontenc}
\usepackage{hyperref}
\title{Makhuwa-Shirima Language Resources for K-12 Curriculum Development}
\author{Compiled by Claude Sonnet 5 (claude-sonnet-5) for Andr\'as Kornai}
\date{}
\begin{document}
\maketitle

\section{Introduction}
Makhuwa-Shirima (ISO 639-3 vmk; Glottocode makh1260; also called Shirima, Xirima, West Makua) is a Bantu language of northern Mozambique with about 500,000 speakers according to the reference description. Searches found scripture and general reference pages. No grammar, corpus, or school materials specific to Makhuwa-Shirima turned up.

\section{Primary Educational Resources}
No primers, textbooks, or school programmes specific to Makhuwa-Shirima were found in search results.

\section{Linguistic and Reference Materials}
\begin{itemize}
\item Wikipedia article on the wider Makhuwa language. It mentions Portuguese-language SIL materials on Central Makhuwa and a survey of Coastal dialects, but these are not specific to Shirima: \url{https://en.wikipedia.org/wiki/Makhuwa_language}
\end{itemize}

\section{Religious and Community Resources}
\begin{itemize}
\item ScriptureEarth page for Makhuwa-Shirima, with Bible text and audio links: \url{https://www.scriptureearth.org/00eng.php?iso=vmk}
\item Joshua Project profile. It reports a New Testament (2000) and a complete Bible (2021): \url{https://joshuaproject.net/people_groups/13422}
\end{itemize}

\section{Digital and Audio Resources}
Audio recordings from the Global Recordings Network are mentioned on the ScriptureEarth page above. A direct URL was not seen in search results.

\section{Curriculum Development Context}
A published Bible gives a written standard and a text base. Orthography, dictionary, and teaching materials were not found online and would need to be located through SIL Mozambique or local education authorities. Sources disagree on where the language is spoken, so confirm this locally.
\end{document}